\subsection{疲劳分析}
\label{sec:fatigue-3in1}

\subsubsection{取值}

“三合一”工具三个模块的安全评定以有限元静力分析为基础，疲劳寿命则须另行计算。本章五类工具采用同一套取值方法，故本节先说明取值的依据，其后各工具只列各自的数值与结果。评定所需确定的量有四项：材料的 S-N 曲线、载荷谱与设计循环数、参与评定的应力值、由应力与循环数计算寿命及安全系数的算式。前三项均依赖具体取值，取值不同则结论不同，故逐项说明如下。

\paragraph{材料强度与 S-N 曲线}
疲劳评定所需的材料参数为抗拉强度、屈服强度与对称循环下的疲劳极限。材料强度取零件图技术要求的下限值。钻具产品的设计以最低保证性能为依据，图纸下限对应可复验的合同指标，而手册所列典型值为统计均值，两者对调质合金钢的差异可达 10\% 以上；若采用典型值，承载能力与疲劳寿命都将被高估。

三种牌号均无可直接引用的 S-N 曲线。4330V 为钻具专用钢号，公开标准中缺乏成组疲劳试验数据；4145 与 4145H 虽有疲劳极限数据，但手册给出的是 $\phi10$ 抛光旋转弯曲试样的结果，与本构件机加工表面、$\phi88\sim\phi105.5$\,mm 的截面尺寸及轴向受载状态均不相符，直接引用将明显高估寿命。因此统一按《机械设计手册》的方法由抗拉强度估算，再用 Marin 五项系数修正至构件的实际状态：
\begin{equation}
\sigma_{-1}=0.5\,\sigma_u\,k_a k_b k_c k_d k_e
\label{eq:marin}
\end{equation}
以毛刷刮管器与液压刮管器承力件的 4145H（825/965 档）为例，五项系数的取值及依据如下。

\begin{itemize}
	\item 表面加工系数 $k_a$。零件图要求机加工面 $Ra1.6$，按手册的机加工曲线 $k_a=4.51\sigma_u^{-0.265}$（$\sigma_u$ 以 MPa 计），代入 965\,MPa 得 0.730。若按抛光试样取值则接近 0.9，将遗漏实际加工面的缺口效应。
	\item 尺寸系数 $k_b$。危险截面直径 $\phi105.5$\,mm，疲劳极限随截面尺寸增大而下降，按手册折减取 0.727。
	\item 载荷类型系数 $k_c$。轴向拉压取 0.85，较旋转弯曲低 15\%。扭转工况不另取系数，因评定采用 von Mises 等效应力，其量纲已折算为当量拉应力。
	\item 温度系数 $k_d$。服役温度不超过 150\,℃，按常温 22\,℃ 取 1.0。
	\item 可靠度系数 $k_e$。按 99\% 可靠度、变异系数 0.08 计，$k_e=1-2.326\times0.08=0.814$。
\end{itemize}

五项系数之积为 $0.730\times0.727\times0.85\times1.0\times0.814=0.367$，代入式~\eqref{eq:marin} 得 $1\times10^{6}$ 次处的疲劳极限 $\sigma_{-1}=0.5\times965\times0.367\approx177.2$\,MPa。

低周端的控制点取 $1\times10^{3}$ 次处的 $0.9\sigma_u=868.5$\,MPa。该值已超过屈服强度 825\,MPa，材料在达到此应力之前即已屈服，曲线不应继续上延，故按屈服强度截断，取 825\,MPa。若照 $0.9\sigma_u$ 取值，等于假定材料能承受高于自身屈服强度的循环载荷。以两个控制点在对数坐标上确定幂律关系，并用 $1\times10^{3}\sim1\times10^{7}$ 次间的五个数值点拟合，得
\begin{equation}
\sigma_a=3\,840\,N^{-0.2227}\qquad(\mathrm{MPa},\ R=-1)
\label{eq:sn-825}
\end{equation}
五个数据点依次为 825、494、296、177 与 106\,MPa，拟合线在 $1\times10^{3}$ 次处算得 825.2\,MPa，与截断值相差 0.02\%，表明截断未破坏曲线的连续性。同一流程用于其余各材料档，所得六条曲线列于表~\ref{tab:sn-params}。

\begin{table}[htbp]
\centering
\caption{本章六条 S-N 曲线的参数（$R=-1$，$1\times10^{3}\sim1\times10^{7}$ 次五点拟合）}
\label{tab:sn-params}
\small
\setlength{\tabcolsep}{4pt}
\begin{tabular}{@{}lccccccc@{}}
\hline
材料档 & $\sigma_s$/MPa & $\sigma_u$/MPa & $k_a k_b k_c k_d k_e$ & $S_e(10^{6})$/MPa & $\sigma_a(10^{3})$/MPa & $C$ & $b$ \\
\hline
4330V-155KSI & 930 & 1\,033 & 0.371 & 191.5 & 929.7 & 4\,512.7 & $-0.2287$ \\
4330V-1068 & 1\,068 & 1\,180 & 0.358 & 211.3 & 1\,062 & 5\,337 & $-0.2337$ \\
4145H-930 & 930 & 1\,080 & 0.367 & 198.1 & 930$^{*}$ & 4\,366 & $-0.2239$ \\
4145H-825 & 825 & 965 & 0.367 & 177.2 & 825$^{*}$ & 3\,840 & $-0.2227$ \\
4145H-760 & 760 & 965 & 0.378 & 182.3 & 760$^{*}$ & 3\,169 & $-0.2067$ \\
4145 & 827 & 965 & 0.378 & 182.3 & 827$^{*}$ & 3\,752 & $-0.2189$ \\
\hline
\end{tabular}

\vspace{2pt}
{\footnotesize 注：$^{*}$ 表示 $0.9\sigma_u$ 已超过屈服强度的档位，该点按 $\sigma_s$ 截断；$S_e$ 为 $1\times10^{6}$ 次处的疲劳极限；$k_a k_b k_c k_d k_e$ 一列为五项系数之积。}
\end{table}

表~\ref{tab:sn-params} 中 4145（827/965）与 4145H（825/965）两档的抗拉强度仅差 2\,MPa，但危险截面直径由 $\phi88$ 增至 $\phi105.5$\,mm，尺寸系数由 0.748 降至 0.727，$1\times10^{6}$ 次处的疲劳极限由 182.3\,MPa 降至 177.2\,MPa，幂指数由 $-0.2189$ 变为 $-0.2227$，同一应力水平下的寿命相差约 10\%。因此同一牌号的不同图纸版本不能共用一条曲线，在有限元软件的材料库中应分别建立条目，不能按牌号名称复用。

\begin{figure}[H]
\centering
\begin{tikzpicture}
\begin{axis}[
	width=0.88\textwidth,
	height=0.55\textwidth,
	xmode=log,
	log basis x=10,
	xmin=900, xmax=1.6e7,
	ymin=90, ymax=900,
	xtick={1000,10000,100000,1000000,10000000},
	xticklabels={$10^{3}$,$10^{4}$,$10^{5}$,$10^{6}$,$10^{7}$},
	ytick={100,200,300,400,500,600,700,800},
	xlabel={应力循环次数 $N$ / 次},
	ylabel={交变应力幅 $\sigma_a$ / MPa},
	label style={font=\small},
	tick label style={font=\small},
	grid=major,
	grid style={densely dotted, gray!45},
	axis line style={thick},
	legend style={font=\tiny, draw=none, fill=white, fill opacity=0.85, text opacity=1, at={(0.98,0.98)}, anchor=north east},
	legend cell align=left,
]
\addplot[very thick, red!75!black, domain=1000:10000000, samples=140] {3840*x^(-0.2227)};
\addlegendentry{拟合曲线 $\sigma_a=3\,840\,N^{-0.2227}$}
\addplot[only marks, mark=*, mark size=1.9pt, black] coordinates {(1000,825) (10000,494) (100000,296) (1000000,177) (10000000,106)};
\addlegendentry{五点时数据}
\addplot[dashed, thick, blue!65!black, domain=1e6:1.25e7] {177.2};
\addlegendentry{$1\times10^{6}$ 次处 $S_e=177.2$\,MPa}
\addplot[only marks, mark=star, mark size=5pt, orange!90!black] coordinates {(9.1e5,180.8)};
\addlegendentry{拉伸工况工作点}
\node[anchor=south east, font=\scriptsize] at (axis cs:8.9e5,192) {工作点 180.8\,MPa（$9.1\times10^{5}$ 次）};
\node[anchor=west, font=\scriptsize] at (axis cs:1.35e3,845) {$1\times10^{3}$ 次处截断值 825\,MPa};
\end{axis}
\end{tikzpicture}
\caption{4145H（825/965 档）的 S-N 曲线。$1\times10^{3}$ 次处 $0.9\sigma_u=868.5$\,MPa 超过屈服强度 825\,MPa，按屈服截断；$1\times10^{6}$ 次处由 $0.5\sigma_u$ 乘 Marin 乘积 0.367 得 177.2\,MPa；五点拟合得 $\sigma_a=3\,840\,N^{-0.2227}$。星号为液压刮管器 400\,kN 拉伸工况的工作点}
\label{fig:sn-curve}
\end{figure}

该曲线有两种用法，本文均有采用：计算寿命时由应力幅反解 $N=(C/\sigma_{a,\mathrm{eq}})^{1/(-b)}$；计算疲劳安全系数时由 $1\times10^{3}$ 次处的许用应力幅正查 $\sigma_e$。图~\ref{fig:sn-curve} 中标出了液压刮管器拉伸工况的工作点，等效对称应力幅 180.8\,MPa 落于 4145H-825 曲线上，对应寿命 $9.1\times10^{5}$ 次，其余工况按同一方式定位。超过 $1\times10^{7}$ 次的部分按水平段延伸，该区域已属无限寿命，损伤增量可忽略。

\paragraph{载荷谱与平均应力修正}
循环特征以应力比 $R=\sigma_{\min}/\sigma_{\max}$ 描述：$R=-1$ 为对称循环，载荷在 $\pm\sigma$ 之间反复；$R=0$ 为脉动循环，载荷由零加至峰值再卸载回零，此时 $\sigma_a=\sigma_m=\sigma_{\max}/2$；$R=1$ 载荷恒定，不产生疲劳损伤。

起下钻的受力过程为"空载—提拉—卸载"，最低载荷为零，故取 $R=0$。抗扭工况分两种情形：上卸扣的扭矩由零加至峰值，仍按 $R=0$；而正反转、倒划眼与憋钻释放会使扭矩换向，大容量强磁工具与多功能过滤工具的抗扭工况按完全反向 $R=-1$ 取定，以覆盖最不利情形。完全反向时平均应力为零、应力幅等于应力峰值，较脉动循环严格一倍，所得寿命相应偏低。

材料的 S-N 曲线在对称循环下测定，直接用于 $R=0$ 的工况将高估寿命，须先将平均应力折算为等效应力幅。本文采用 Goodman 准则：
\begin{equation}
\frac{\sigma_a}{\sigma_{-1}}+\frac{\sigma_m}{\sigma_u}=1
\quad\Longrightarrow\quad
\sigma_{a,\mathrm{eq}}=\frac{\sigma_a}{1-\sigma_m/\sigma_u}
\label{eq:goodman}
\end{equation}
选择 Goodman 而不用 Gerber 或 Soderberg，原因有三：Gerber 抛物线对高强度钢偏于乐观，用于钻具将低估风险；Soderberg 以屈服强度为界，过于保守；Goodman 为 ASME BPVC VIII-2 与 API 各规范默认采用的准则，与静强度判据同源，两者的结果可并列于同一表中比较。

设计循环数按单套工具服役 200 趟起下钻计，$n=200$ 次。合同与图纸均未规定使用趟数，该值属设计假定，因此本文同时给出计算寿命 $N$：累计使用不超过 $N$ 次时累积损伤恒小于 1.0，结论并不依赖 200 这一假定。若按每趟 5 次行程计，$n=1\,000$ 次，各工况损伤按比例放大五倍，三合一工具三个模块中最严的一档为 $4.4\times10^{-3}$，距判据 1.0 仍有两个量级以上。

判定采用三项指标：静强度安全系数 $n_s=\sigma_s/\sigma_{\max}\ge1.5$；200 次累积损伤 $D=\sum n_i/N_i\le1.0$，按 Miner 线性累积准则计算；疲劳安全系数 $n_f=1/(\sigma_a/\sigma_e+\sigma_m/\sigma_u)\ge1.5$，其中 $\sigma_e$ 为该材料 $1\times10^{3}$ 次处的许用应力幅。$n_f$ 的物理意义为载荷尚可放大的倍数，其同时放大了应力幅与平均应力，与 Goodman 图上的射线长度相对应，因而较单一损伤指标更能反映安全余量。

\paragraph{评定应力的确定}
\label{ss:fatigue-probe}

有限元云图的峰值不能直接作为评定应力。模型中的几何尖角（螺纹根部、台阶过渡）、接触终止线与固定约束端面，在弹性解下应力随网格加密单调增大而不收敛，峰值仅反映网格尺度，并不反映材料的真实受力。以液压刮管器 400\,kN 拉伸工况为例：云图峰值 484.02\,MPa 位于连接体与本体配合端面的接触终止线，同一截面距峰值 3$\sim$6\,mm 处的应力为 304.51\,MPa，两者相差 37\%。若以峰值评定，静强度安全系数降为 1.70；以探针值评定则为 2.71。同一工况两种口径相差一倍有余，表明取样位置本身即是结论的组成部分，须先确立取数规则。

取评定应力的规则为：在峰值附近 3$\sim$6\,mm、应力梯度已趋平缓处取三个节点，要求位于同一环带、网格尺寸相近，且三点之间的相对差小于 5\%，以确认已脱离奇异区；取三点中的最大值作为评定应力。液压刮管器拉伸工况所取的 301.52、206.03 与 304.51\,MPa 三个节点中，同环带的 301.52 与 304.51\,MPa 相差 1.0\%，表明该处梯度已趋平缓；206.03\,MPa 位于另一侧过渡区，用于观察应力沿轴向的衰减。评定取 304.51\,MPa。图~\ref{fig:probe} 给出了峰值、取样带与评定值三者的关系。

\begin{figure}[H]
\centering
\begin{tikzpicture}
\begin{axis}[
	width=0.88\textwidth,
	height=0.45\textwidth,
	xmin=-0.6, xmax=10.6,
	ymin=170, ymax=600,
	xtick={0,2,4,6,8,10},
	ytick={200,250,300,350,400,450,500,550},
	xlabel={距奇异点的轴向距离 / mm},
	ylabel={等效应力 $\sigma$ / MPa},
	label style={font=\small},
	tick label style={font=\small},
	grid=major,
	grid style={densely dotted, gray!45},
	axis line style={thick},
]
\fill[green!12] (axis cs:3,170) rectangle (axis cs:6,600);
\addplot[smooth, tension=0.72, very thick, blue!65!black] coordinates {(0,484.02) (0.8,425) (1.6,368) (2.4,330) (3,304.51) (4.5,301.52) (6,290) (7.5,255) (9,222) (10,206.03)};
\addplot[dashed, thick, red!70!black, domain=0:10.6] {304.51};
\addplot[only marks, mark=*, mark size=2.4pt, red!80!black] coordinates {(0,484.02) (3,304.51) (4.5,301.52) (10,206.03)};
\node[anchor=south west, font=\scriptsize, align=left] at (axis cs:0.3,495) {峰值 484.02\,MPa\\随网格加密继续增大};
\node[anchor=south, font=\scriptsize] at (axis cs:4.5,545) {评定取样带 3$\sim$6\,mm};
\node[anchor=north, font=\scriptsize, align=left] at (axis cs:3.2,300) {304.51 与 301.52\\相差 1.0\%};
\node[anchor=north west, font=\scriptsize] at (axis cs:10.4,204) {206.03};
\node[anchor=east, font=\scriptsize] at (axis cs:9.9,335) {评定应力 304.5\,MPa};
\end{axis}
\end{tikzpicture}
\caption{评定应力的取值规则（示意，数值取液压刮管器 400\,kN 拉伸工况）。峰值位于接触终止线，随网格加密继续增大，不参与评定；距峰值 3$\sim$6\,mm 的取样带内 301.52\,MPa 与 304.51\,MPa 相差 1.0\%，取大值 304.51\,MPa 作为评定应力；206.03\,MPa 为沿轴向更远处的衰减点}
\label{fig:probe}
\end{figure}

三类伪应力的来源不同，处理方式相同而机理各异，应分别说明。几何尖角引起的奇异（如可变径刮管器内孔台阶根部，峰值 2\,156.6\,MPa）源于网格离散，加密网格只会使峰值继续增大；接触终止线处的高应力（如液压刮管器配合端面）源于薄壁环件贴合面的边界条件；固定约束端面的伪应力（如毛刷刮管器拉伸工况的 638.49\,MPa）则因约束强制端面保持平面，与邻近材料的自由变形趋势相冲突，与零件几何无关。后者可用应力集中系数反证：毛刷刮管器本体名义应力为 107.7\,MPa，反算峰值处的集中系数为 5.93，而机加工台阶、倒角与退刀槽的集中系数通常在 1.5$\sim$3.6 之间；同一模型上台肩处的集中系数为 3.64，恰落于合理区间。同一模型中两类特征同时出现，说明前者源于边界条件、后者源于几何；若将两者均视为"应力集中"，则无法据此提出改进方向。

峰值虽不作判据，仍予单列备查，因其反映该处的梯度陡峭程度，并提示结果对取样位置较为敏感。液压刮管器扭转工况的峰值为 799.14\,MPa，已达屈服强度的 97\%，以峰值口径计算的静强度安全系数仅为 1.03；按规则该值不参与评定，但其表明该配合端面是扭转工况的薄弱环节，后续应加密网格复核，而不能将 1.03 作为结论。

三张云图同源，取数时应相互印证：寿命图的最小值、损伤图的最大值与安全系数图的最小值须落在同一节点，若分处三地，则说明材料指派、载荷步或循环类型设置有误，应先改正设置再读数。损伤云图的读数还须注明设计循环数基准：本文部分工况按 $N_d=1$ 次输出，读数即单次损伤；其余按 $N_d=1\times10^{3}$ 次输出，读数须除以 1\,000 再乘 200 方能与 200 趟对齐。云图最小值与探针值允许 5\%$\sim$8\% 的差异，来源有三：最小值可能落于奇异点上，与探针不在同一节点；云图按单元积分点插值至节点，探针读取的是节点解；S-N 曲线在 $1\times10^{3}$ 次处按屈服截断，幂律在该点附近并不严格通过截断值。差值在此范围内视为同源，超过 10\% 则表明探针落于奇异点或陡梯度区，应重新布点。

\paragraph{有限元疲劳模块的设定}
\label{ss:fatigue-module}
本章各工具的寿命、损伤与安全系数云图均由 ANSYS 的 Fatigue Tool 按应力寿命（S-N）法给出，各项设定在全部工具中保持一致：分析类型取“应力寿命”；载荷类型按工况的循环特征选取——拉伸与脉动受扭取“基于零”（Zero-Based），载荷在零与峰值之间往复，对应 $R=0$，软件由此取 $\sigma_a=\sigma_m=\sigma_{\max}/2$；换向的抗扭工况取“完全反向”（Fully Reversed），对应 $R=-1$，$\sigma_a=\sigma_{\max}$、$\sigma_m=0$；平均应力理论取 Goodman，即软件先按 $\sigma_{a,\mathrm{eq}}=\sigma_a/(1-\sigma_m/\sigma_u)$ 折算到等效对称应力幅，再以该值查材料 S-N 曲线；材料疲劳强度因子 $K_f=1$，不另加缺口折减，螺纹根部、退刀槽等缺口效应由模型几何本身体现；应力分量取带符号的 von Mises，使拉压与扭转在同一量纲上与 S-N 曲线对应；设计循环数取 $N_d=1\times10^{3}$ 次，损伤云图的读数即 $N_d/N_i$（个别工况按 $N_d=1$ 次输出，读数即单次损伤，引用时须先核对基准）。由于 Goodman 修正与表~\ref{tab:fatigue-chain} 的换算口径一致，云图给出的寿命、损伤与安全系数可直接与表中数值对照，二者允许 5\%$\sim$8\% 的差异（见上一段）；对照时须先确认载荷类型与循环特征相符，若某工况实际为完全反向而软件按基于零设置，读数会偏保守约一个 Goodman 因子，不能直接引用。

\paragraph{由应力计算寿命与安全系数}
\label{ss:fatigue-chain}
取值确定之后，各项结论按固定顺序计算，每一步只对应一个算式，如表~\ref{tab:fatigue-chain} 所列。全部算式的输入仅有评定应力 $\sigma_{\max}$、材料的 $\sigma_s$ 与 $\sigma_u$、以及 S-N 曲线的 $C$ 与 $b$。更换材料或载荷时只改这三组输入，算式与顺序不变，各工具的结果因而可以横向比较。

\begin{table}[htbp]
\centering
\caption{由评定应力到三项判据的计算流程}
\label{tab:fatigue-chain}
\small
\begin{tabularx}{\textwidth}{@{}lXc@{}}
\hline
计算量 & 算式与依据 & 液压刮管器拉伸算例 \\
\hline
应力幅与平均应力 & $R=0$：$\sigma_a=\sigma_m=\sigma_{\max}/2$；起下钻载荷在零与峰值之间往复 & 152.26\,MPa \\
等效对称应力幅 & $\sigma_{a,\mathrm{eq}}=\sigma_a/(1-\sigma_m/\sigma_u)$，Goodman 修正 & 180.8\,MPa \\
计算寿命 & $N=(C/\sigma_{a,\mathrm{eq}})^{1/(-b)}$，S-N 幂律反解 & $9.1\times10^{5}$ 次 \\
200 次累积损伤 & $D=n/N$，$n=200$，Miner 线性累积 & $2.2\times10^{-4}$ \\
疲劳安全系数 & $n_f=1/(\sigma_a/\sigma_e+\sigma_m/\sigma_u)$，$\sigma_e=825$\,MPa & 2.92 \\
静强度安全系数 & $n_s=\sigma_s/\sigma_{\max}$ & 2.71 \\
\hline
\end{tabularx}
\end{table}

以液压刮管器 400\,kN 拉伸工况为例说明计算过程。评定应力 $\sigma_{\max}=304.51$\,MPa，按 $R=0$ 得 $\sigma_a=\sigma_m=152.26$\,MPa；经 Goodman 折算得 $\sigma_{a,\mathrm{eq}}=152.26/(1-152.26/965)=180.8$\,MPa；由 4145H-825 的幂律反解得 $N=(3\,840/180.8)^{1/0.2227}=9.1\times10^{5}$ 次；200 趟的累积损伤 $D=200/(9.1\times10^{5})=2.2\times10^{-4}$；疲劳安全系数 $n_f=1/(152.26/825+152.26/965)=2.92$；静强度安全系数 $n_s=825/304.51=2.71$。三项判据分别为 $2.71\ge1.5$、$2.2\times10^{-4}\le1.0$、$2.92\ge1.5$，均满足要求，寿命裕度为 $9.1\times10^{5}/200\approx4\,550$ 倍。

上述结果可与软件输出逐项印证，这是取数可信的依据。该工况寿命云图的最小值为 67\,523 次，与按峰值 484.02\,MPa 手算的 $6.72\times10^{4}$ 次相差 0.4\%；损伤云图读数为 $1.48\times10^{-2}$（按 1\,000 次循环），与手算的 $1.49\times10^{-2}$ 一致；节点 41108 的安全系数读数为 2.9212，与按 304.51\,MPa 算得的 2.92 相同。三处互证表明载荷、边界条件与材料曲线的输入与手算同源，云图读数可用于评定；若三者不符，应先核查载荷步、材料指派与循环类型，而不宜先怀疑计算结果。

\subsubsection{评定}

三合一工具三个模块的额定载荷同为 400\,kN 与 8\,kN$\cdot$m，材料分两档：脉冲短节的心轴与上下接头为 4145（827/965），毛刷刮管器的芯轴与下接头、液压刮管器的四个承力件为 4145H（825/965）。两档抗拉强度仅差 2\,MPa，但危险截面不同，尺寸系数由 0.748 降为 0.727，疲劳极限由 182.3\,MPa 降为 177.2\,MPa，故各建一条曲线：脉冲短节用 $\sigma_a=3\,752\,N^{-0.2189}$，毛刷刮管器与液压刮管器用 $\sigma_a=3\,840\,N^{-0.2227}$，$1\times10^{3}$ 次处的许用应力幅分别为 827\,MPa 与 825\,MPa。

拉伸工况的评定应力取自各模块的危险截面：脉冲短节 340.09\,MPa（节点 22369，距奇异点 4$\sim$6\,mm，相邻节点 319.11\,MPa），毛刷刮管器 392.26\,MPa（节点 17294，相邻 370.49\,MPa），液压刮管器 304.51\,MPa（节点 41108，同环带 301.52\,MPa）。扭转工况按 8\,kN$\cdot$m 施加，评定应力分别为 464.47、384.16 与 376.16\,MPa，其中毛刷刮管器取台肩环带值而非棱边值。两工况均按 $R=0$、200 趟设计循环数评定，结果列于表~\ref{tab:fatigue_3in1}。

\begin{table}[htbp]
\centering
\caption{三合一工具三个模块的疲劳评定结果（拉伸 400\,kN、扭转 8\,kN$\cdot$m，$R=0$，200 趟）}
\label{tab:fatigue_3in1}
\begin{tabularx}{\textwidth}{lCCC}
\hline
项目 & 脉冲短节 & 毛刷刮管器 & 液压刮管器 \\
\hline
材料档 & 4145（827/965） & 4145H（825/965） & 4145H（825/965） \\
承力零件数 & 3 & 2 & 4 \\
S-N 曲线 $\sigma_a=CN^b$ & $3\,752\,N^{-0.2189}$ & $3\,840\,N^{-0.2227}$ & $3\,840\,N^{-0.2227}$ \\
拉伸评定应力 / MPa & 340.1 & 392.3 & 304.5 \\
应力幅与平均应力 / MPa & 170.1 & 196.1 & 152.3 \\
等效对称应力幅 / MPa & 206.4 & 246.2 & 180.8 \\
计算寿命 / 次 & $5.7\times10^{5}$ & $2.3\times10^{5}$ & $9.1\times10^{5}$ \\
200 次累积损伤 & $3.4\times10^{-4}$ & $8.8\times10^{-4}$ & $2.2\times10^{-4}$ \\
疲劳安全系数 & 2.62 & 2.27 & 2.92 \\
静强度安全系数 & 2.43 & 2.10 & 2.71 \\
\hline
抗扭评定应力 / MPa & 464.5 & 384.2 & 376.2 \\
抗扭等效对称应力幅 / MPa & 305.6 & 239.8 & 233.6 \\
抗扭计算寿命 / 次 & $9.4\times10^{4}$ & $2.6\times10^{5}$ & $2.9\times10^{5}$ \\
抗扭 200 次累积损伤 & $2.1\times10^{-3}$ & $7.8\times10^{-4}$ & $7.0\times10^{-4}$ \\
抗扭疲劳安全系数 & 1.92 & 2.32 & 2.36 \\
抗扭静强度安全系数 & 1.78 & 2.15 & 2.19 \\
\hline
\end{tabularx}
\end{table}

表~\ref{tab:fatigue_3in1} 中每一行均可由本节的计算流程逐列复核。以毛刷刮管器拉伸工况为例：$\sigma_a=\sigma_m=196.13$\,MPa，$\sigma_{a,\mathrm{eq}}=196.13/(1-196.13/965)=246.2$\,MPa，$N=(3\,840/246.2)^{1/0.2227}=2.3\times10^{5}$ 次，$D=200/(2.3\times10^{5})=8.8\times10^{-4}$，$n_f=1/(196.13/825+196.13/965)=2.27$，$n_s=825/392.26=2.10$。静强度安全系数 $n_s$ 与疲劳安全系数 $n_f$ 数值接近而含义不同：前者仅为峰值应力与屈服强度之比，后者将应力幅与平均应力同时放大至许用边界，故 $n_f$ 恒不小于 $n_s$。

三个模块的拉伸评定应力在 304$\sim$392\,MPa 之间，计算寿命为 $2.3\times10^{5}\sim9.1\times10^{5}$ 次，200 趟累积损伤均在 $10^{-3}$ 量级以下，寿命裕度 1\,150$\sim$4\,550 倍，全部满足判据。与 1\,000\,kN 级工具（评定应力 588$\sim$659\,MPa、寿命 $3.0\times10^{4}\sim3.7\times10^{4}$ 次）相比，三个模块的裕度明显更宽，原因在于承力截面按 400\,kN 设计、壁厚相对充裕，而载荷仅为其 40\%。扭转工况的评定应力为 376$\sim$465\,MPa，损伤为 $7.0\times10^{-4}\sim2.1\times10^{-3}$，安全系数为 1.78$\sim$2.36，同样满足要求，并不构成较拉伸更严酷的控制工况。

有两处细节表明"取哪一点的应力"确实会改变结论。其一，毛刷刮管器拉伸工况的峰值 638.49\,MPa 落于固定约束端面，并不位于任何几何特征上：按名义应力 107.7\,MPa 反算，该处应力集中系数为 5.93，远超机加工台阶与退刀槽的合理区间，而同一模型上台肩处的 3.64 才源于几何。约束端面的伪应力可用两种方法消除：将固定支撑改为六自由度远程位移（行为取 Deformable），或在两端各延长 50$\sim$100\,mm 过渡几何后再施加约束；台肩处的评定值基本不变，此点已由相邻节点的采样一致性佐证。

其二，毛刷刮管器扭转工况的棱边读数为 1005.4、722.16 与 710.35\,MPa，相邻节点相差 41\%；将该处棱边网格重划后，同一位置的读数由 761.32 降至 700.68\,MPa，降幅 8\%，说明该处应力水平由尖角几何与网格离散控制，属数值奇异。离开棱边、位于台肩环带上的 384.16、381.34 与 363.72\,MPa 仅相差 5.4\%，才是可用于工程的评定值，本文取其中的 384.16\,MPa。

各模块的控制部位与相应措施如下：脉冲短节为心轴台阶根部与螺纹退刀槽终止线，应保持不小于 R3 的过渡圆角，退刀槽底部按 R1 以上加工；毛刷刮管器为台肩过渡处，与图纸"抛光圆角不超过 R0.2"的去应力集中措施一致，该圆角在建模与网格划分中必须保留，简化处理将使该处应力被抹平、寿命被高估；液压刮管器为连接体与本体配合端面的环形过渡带，建议将该处圆角加工至 R2 以上，并对配合端面加密一档网格作无关性复核，其扭转工况的峰值口径静强度安全系数仅为 1.03，虽不参与评定，仍表明该环带对扭矩较为敏感。三个模块的成品均应按 NB/T\,47013.4 进行磁粉检测。
